\documentclass[11pt]{article}
\usepackage[a4paper, margin=1in]{geometry}
\usepackage{amsmath, amssymb, amsthm, bm}
\usepackage{xcolor}
\usepackage{listings}
\usepackage{hyperref}
\usepackage{tikz}
\usepackage{booktabs}
\usepackage{enumitem}

\hypersetup{colorlinks=true, urlcolor=blue, linkcolor=blue}

% -- Csurös-style notation (CM09 SI / Csurös 2021 arXiv) ------------------
\newcommand{\R}{\mathbb{R}}
\newcommand{\N}{\mathbb{N}}
\newcommand{\Z}{\mathbb{Z}}
\newcommand{\Tree}{\mathcal{T}}
\newcommand{\Leaves}{\mathcal{L}}
\newcommand{\pa}{\mathrm{pa}}
\newcommand{\ch}{\mathrm{ch}}
\newcommand{\troot}{R}
\newcommand{\Prof}{\boldsymbol{\xi}}   % per-family copy-number profile
\newcommand{\rate}{\kappa}             % gain shape
\newcommand{\dup}{\lambda}             % duplication rate
\newcommand{\loss}{\mu}                % loss rate
\newcommand{\edgelen}{t}               % edge length
\newcommand{\survp}{\tilde p}          % survival probability (CM09 SI)
\newcommand{\survq}{\tilde q}          % Pólya per-event parameter
\newcommand{\Ll}{\mathcal{L}}          % log-likelihood
\DeclareMathOperator{\Var}{Var}
\DeclareMathOperator{\Cov}{Cov}
\DeclareMathOperator{\logit}{logit}

\newtheorem{theorem}{Theorem}
\newtheorem{proposition}[theorem]{Proposition}
\newtheorem{lemma}[theorem]{Lemma}
\newtheorem{corollary}[theorem]{Corollary}
\theoremstyle{definition}
\newtheorem{definition}[theorem]{Definition}
\newtheorem{remark}[theorem]{Remark}
\newtheorem{example}[theorem]{Example}

\title{Autocorrelated rate priors on phylogenies:\\
       a Thorne--Kishino--Painter relaxed clock for the\\
       gain--loss--duplication birth--death model}
\author{\texttt{recount} project\\
        \textit{notation following CM09 SI and Cs\H{u}r\"os 2021 arXiv}}
\date{\today}

\begin{document}
\maketitle

\begin{center}
\fbox{\parbox{0.9\textwidth}{\small\textbf{Provenance and disclaimer.}
This document was drafted by Claude (Anthropic's LLM, model
\texttt{claude-opus-4-7[1m]}) in collaboration with the human author,
as part of an interactive Claude Code session
(\url{https://claude.com/claude-code}). The Thorne--Kishino--Painter
autocorrelated-rate prior is standard (Thorne, Kishino \& Painter,
\emph{MBE} 15:1647--1657, 1998); the per-edge gradient was
finite-difference validated to relative error $\le 10^{-6}$ against the
analytical implementation in
\texttt{native/src/recount\_brownian\_prior.c} and the Python reference
in \texttt{validation/\_shared.py:\_brownian\_prior\_and\_grad\_python}.
Treat as ``LLM-drafted, human-checked, FD-validated'' engineering
documentation rather than peer-reviewed manuscript content.}}
\end{center}

\begin{abstract}
We equip the \texttt{recount} maximum-a-posteriori fitter with a
tree-structured prior that couples the rates of adjacent branches.
The construction adapts the relaxed-clock prior of Thorne, Kishino
\& Painter (\emph{MBE} 1998), placing a Brownian motion on the
log-rate process as it walks down the phylogeny.  We give the
explicit log-density (Eq.~\eqref{eq:logprior-final}) and gradient
(Eq.~\eqref{eq:grad-edge}), prove the prior is proper
(Theorem~\ref{thm:proper}) and that the resulting MAP objective is
coercive (Theorem~\ref{thm:coercive}). The implementation is
$\mathcal{O}(N)$ per evaluation and matches the analytical native C
kernel to FD relative error~$\le10^{-6}$.
\end{abstract}

% =================================================================
\section{Biological motivation}
\label{sec:motivation}
% =================================================================

A gain--loss--duplication (GLD) phylogenetic fit estimates four
quantities per branch of the species tree: a gain shape $\rate_v$
controlling \emph{de novo} acquisition of new families, a duplication
rate $\dup_v$, a loss rate $\loss_v$ (conventionally fixed at one to
set the timescale), and a branch length $\edgelen_v$. On a tree with
$N$ nodes that gives $\sim 3N$ free parameters. For the archaeal
arc269 dataset ($N=537$) this is $\sim 1{,}600$ rate parameters.

In every published reconstruction we are aware of these rates are
treated as completely independent across branches. From an
evolutionary-biology standpoint that is a strong, and usually
unwarranted, assumption: the descendants of a recent common ancestor
share an ecological niche, a mutational milieu, an effective population
size, and (often) a metabolic strategy --- all of which influence the
rate of gene gain, loss, and duplication. A sudden, large jump in
inferred gain rate from one branch to its child is biologically
implausible most of the time, but the independent prior treats it as
just as likely as no jump.

What is required is a prior that says \emph{``adjacent branches should
have similar rates, with gradually growing divergence over evolutionary
time.''} This is exactly the structure introduced by Thorne, Kishino
\& Painter (1998, henceforth \textbf{TKP}) for molecular substitution
rates: the log-rate is allowed to wander Brownian-style as one walks
from the root down the tree, with variance growing with branch length.
TKP's prior was an early example of what is now called a
\emph{relaxed clock}: it relaxes the strict-clock assumption (all
branches share one rate) without collapsing to the unstructured
independent-rate limit.

\paragraph{What this document does.} We adapt TKP to GLD rates.
Each axis (gain, duplication, length) gets its own independent
Brownian process on the log-rate as it walks from the root through
the tree. We work in the coordinates used by the optimiser
(\S\ref{sec:setup}), give the explicit log-prior density
(\S\ref{sec:density}) and gradient (\S\ref{sec:gradient}), prove that
the resulting joint prior on rates is a full-rank Gaussian Markov
random field, hence proper (\S\ref{sec:propriety}), and prove that
when combined with the GLD likelihood the MAP objective is coercive
(\S\ref{sec:coercivity}). This last point is the key practical payoff:
it rules out the rate-boundary asymptote analysed in the companion
note \texttt{boundary\_identifiability.tex}.

% =================================================================
\section{Setup and notation}
\label{sec:setup}
% =================================================================

We follow the notation of Cs\H{u}r\"os 2021 (arXiv:2107.11440), with
the survival-parameter convention of the CM09 supplement.

\paragraph{Tree.}
Let $\Tree = (V, E, \troot)$ be a rooted tree with node set $V$ of
cardinality $N$, root $\troot \in V$, leaf set $\Leaves \subset V$,
and edges directed from $\troot$ outward.  For non-root $v$ we write
$\pa(v) \in V$ for its parent. Each edge $\pa(v)\to v$ is identified
with its child endpoint $v$.

\paragraph{Per-edge rates.}
On the incoming edge of each non-root node $v$ we attach four
parameters:
\begin{equation}
\rate_v \;\ge\; 0,\qquad
\dup_v \;\ge\; 0,\qquad
\loss_v \;>\; 0,\qquad
\edgelen_v \;\in\; (0, \infty).
\label{eq:rates}
\end{equation}
$\rate_v$ is the immigration (Pólya shape) parameter for new gene
families landing on the edge, $\dup_v$ the per-lineage duplication
rate, $\loss_v$ the per-lineage loss rate, and $\edgelen_v$ the time
elapsed on the edge. Following CM09 we set $\loss_v \equiv 1$ (it
absorbs the timescale), so the free per-edge parameters are
$(\rate_v, \dup_v, \edgelen_v)$. The root $\troot$ carries a gain
parameter $\rate_\troot$ but no edge length.

\paragraph{The optimiser's coordinates.}
Because rates are positive, the optimiser works on log-coordinates;
for duplication we additionally apply the sub-critical Yule logit
when activated. With $N$ nodes (including root) define the packed
vector
\begin{equation}
x \;=\; \bigl(\underbrace{\log\rate_v}_{v\in V},\;
            \underbrace{\theta_{\dup,v}}_{v\neq\troot},\;
            \underbrace{\log\edgelen_v}_{v\neq\troot}\bigr)
\;\in\; \R^{3N-2},
\label{eq:packing}
\end{equation}
where
\begin{equation}
\theta_{\dup,v} \;=\;
\begin{cases}
\logit\bigl(\dup_v \,/\, \mathrm{Q}\bigr) & \text{(sub-critical Yule, $\mathrm{Q}=1-2^{-30}$)}\\
\log\dup_v & \text{(unconstrained mode)}
\end{cases}
\label{eq:theta-dup}
\end{equation}
matches the Cs\H{u}r\"os Java convention
(\texttt{is\_duprate\_bounded=true}). The gain axis includes the root;
the duplication and length axes have $N-1$ entries each (no root). To
unify the three axes for the prior, write $x^{(a)}_v$ for the entry
of $x$ at node $v$ on axis $a\in\{\rate, \dup, \edgelen\}$, with the
convention that $x^{(\dup)}_\troot$ and $x^{(\edgelen)}_\troot$ are
identified with fixed reference centres (see \S\ref{sec:density}
below).

\paragraph{Data.}
A gene-family profile is a vector
$\Prof \in \Z_{\ge 0}^\Leaves$ of integer copy counts at each leaf.
The data is $\{\Prof_f\}_{f=1}^F$. The log-likelihood
$\Ll(x) := \sum_f \log P(\Prof_f \mid x)$ is computed by the
inside--outside recursion of CM09; this document does not modify it.

% =================================================================
\section{The prior: classical TKP and our v1 simplification}
\label{sec:prior}
% =================================================================

\subsection{Classical TKP on log-rates}

In its original form, TKP places a Gaussian on the log-rate
\emph{increment} between parent and child, with variance proportional
to elapsed evolutionary time:
\begin{equation}
\log r_v - \log r_{\pa(v)}
\;\sim\;
\mathcal{N}\!\bigl(0,\, \sigma_{\mathrm{TKP}}^2 \cdot \edgelen_v\bigr).
\label{eq:tkp-classic}
\end{equation}
This says: the log-rate process is a Brownian motion run along the
tree, with diffusion coefficient $\sigma_{\mathrm{TKP}}^2$. Two
sister branches with short edges between them and a recent common
ancestor will have similar rates; two distant branches separated by
long edges will be free to diverge.

\subsection{Our v1: time-uniform Brownian}

We adopt the structure of \eqref{eq:tkp-classic} with one
simplification: drop the per-edge $\edgelen_v$ weighting, so that
every edge incurs the same variance $\sigma_a^2$ on each rate axis
$a$. Formally,
\begin{equation}
\Delta_a(v)
\;:=\;
x^{(a)}_v - x^{(a)}_{\pa(v)}
\;\overset{\text{iid}}{\sim}\;
\mathcal{N}\!\bigl(0,\, \sigma_a^2\bigr)
\qquad
\text{for every } v\neq \troot,\quad a\in\{\rate,\dup,\edgelen\}.
\label{eq:tkp-v1}
\end{equation}
This is mathematically a Gaussian process on the tree with covariance
function $\Cov(x_u, x_v) = \sigma_a^2 \cdot d(u,v)$ where $d(u,v)$ is
the \emph{topological} (edge-count) distance to the most recent
common ancestor, rather than the metric distance. Biologically it
treats every branch as the same ``unit of evolutionary opportunity''
to revise the rate, irrespective of its $\edgelen_v$.

\begin{remark}[Why drop the $\edgelen_v$ weight?]\label{rem:why-tu}
The classical TKP weight $\edgelen_v$ couples the prior to the very
quantity ($\edgelen_v$ itself) that the optimiser is solving for,
introducing a non-trivial extra chain-rule term whenever
$\edgelen_v$ moves.  The time-uniform variant produces an
$\mathcal{O}(N)$ kernel with no $\edgelen$-coupling, and empirically
captures the same first-order effect (shrinkage between adjacent
branches). The classical form is left for a v2 implementation; see
\S\ref{sec:future}.
\end{remark}

\subsection{Anchoring the prior at the root}

A Brownian motion has no global scale: shifting every $x^{(a)}_v$ by
a common constant leaves all the increments unchanged.  To pin the
absolute scale we add a wide independent Gaussian \emph{root anchor}
\begin{equation}
\log\rate_\troot \;\sim\; \mathcal{N}(\loss^{(\rate)}_\troot,\,
\sigma_\troot^2),
\label{eq:root-anchor-gain}
\end{equation}
where by convention we use $\loss^{(a)}_\troot$ for the prior centre
on axis $a$.  On the duplication and length axes the root does not
have an optimisable rate (the root has no incoming edge), so the
root values are fixed at the prior centres
$\loss^{(\dup)}_\troot, \loss^{(\edgelen)}_\troot$ and
``$\Delta_a(v)$ for $\pa(v)=\troot$'' is interpreted with these
fixed values plugged in.

% =================================================================
\section{Log-density}
\label{sec:density}
% =================================================================

Combining \eqref{eq:tkp-v1} and \eqref{eq:root-anchor-gain}, the
joint log-prior is
\begin{equation}
\boxed{\;
\log p_{\mathrm{B}}(x)
\;=\;
- \tfrac{1}{2}\sum_{a\in\{\rate,\dup,\edgelen\}}\;
   \sum_{v\neq\troot}
      \frac{\Delta_a(v)^2}{\sigma_a^2}
\;-\;
\tfrac{1}{2}\frac{\bigl(\log\rate_\troot - \loss^{(\rate)}_\troot\bigr)^2}
                 {\sigma_\troot^2}
\;+\; \mathrm{const}.
\;}
\label{eq:logprior-final}
\end{equation}
The constant absorbs the normalising factors of the Gaussians; under
the v1 time-uniform variance they do not depend on $x$.

\paragraph{Worth pausing on.}
Equation~\eqref{eq:logprior-final} is just a sum of squared
log-rate increments along the tree --- one per edge, per axis, plus
one root-anchor term on the gain axis. There are no off-tree
interactions, no quadratic couplings between non-adjacent branches.
This is what makes the gradient and the inference $\mathcal{O}(N)$.

% =================================================================
\section{Gradient}
\label{sec:gradient}
% =================================================================

The gradient is obtained by direct differentiation of
\eqref{eq:logprior-final}. Each edge contribution
$-\Delta_a(v)^2/(2\sigma_a^2)$ depends on the child entry $x^{(a)}_v$
and on the parent entry $x^{(a)}_{\pa(v)}$ (the latter only when
$\pa(v)$ is itself not the root, or on the gain axis where the root
is in $x$). The closed form is
\begin{align}
\frac{\partial \log p_{\mathrm{B}}}{\partial x^{(a)}_v}
&= -\frac{\Delta_a(v)}{\sigma_a^2}
&&\text{(contribution from } v \text{ as a child)},
\label{eq:grad-edge}\\
\frac{\partial \log p_{\mathrm{B}}}{\partial x^{(a)}_{\pa(v)}}
&\mathrel{+}= +\frac{\Delta_a(v)}{\sigma_a^2}
&&\text{(contribution from } \pa(v) \text{ as a parent of } v\text{)},
\label{eq:grad-parent}\\
\frac{\partial \log p_{\mathrm{B}}}{\partial \log\rate_\troot}
&\mathrel{+}= -\frac{\log\rate_\troot - \loss^{(\rate)}_\troot}{\sigma_\troot^2}
&&\text{(root anchor)}.
\label{eq:grad-root}
\end{align}
The ``$\mathrel{+}=$'' notation captures the scatter-add semantics:
the parent of any internal node $u$ accumulates contributions from
all of $u$'s children $v\in\ch(u)$. The implementation is a single
pass over the $N-1$ tree edges with $\mathcal{O}(1)$ work per edge
per axis.

\begin{proposition}[Computational complexity]\label{prop:complexity}
The log-prior \eqref{eq:logprior-final} and its full gradient
\eqref{eq:grad-edge}--\eqref{eq:grad-root} cost
$\mathcal{O}(N)$ floating-point operations.
\end{proposition}
\begin{proof}
One pass over the $N-1$ tree edges with three scalar additions per
edge per axis, plus a single root-anchor term. $\square$
\end{proof}

% =================================================================
\section{Theoretical properties}
% =================================================================

\subsection{Propriety: the prior is a full-rank GMRF}
\label{sec:propriety}

We show that the joint prior on the packed vector $x$ defined by
\eqref{eq:logprior-final} is a non-degenerate (full-rank) Gaussian
distribution. Equivalently, it integrates to a finite constant, so
it is a \emph{proper} prior in the Bayesian sense.

\begin{theorem}[Propriety]\label{thm:proper}
Fix $\sigma_a > 0$ for $a \in \{\rate, \dup, \edgelen\}$ and
$\sigma_\troot > 0$. Under the conventions of
\S\ref{sec:setup}--\S\ref{sec:density}, the prior density
\eqref{eq:logprior-final} arises from a Gaussian distribution on
$x \in \R^{3N-2}$ with positive-definite precision matrix
\begin{equation}
\Sigma^{-1} \;=\; \mathrm{blkdiag}\bigl(\Lambda_\rate,\;
                                       \Lambda_\dup,\;
                                       \Lambda_\edgelen\bigr),
\label{eq:precision-block}
\end{equation}
where on the gain axis $\Lambda_\rate = (1/\sigma_\rate^2)\,L_\Tree
+ (1/\sigma_\troot^2)\,e_\troot e_\troot^\top$ with $L_\Tree$ the
unsigned graph Laplacian of $\Tree$, and on the duplication and
length axes $\Lambda_a = (1/\sigma_a^2)\,L^\partial_\Tree$ is the
\emph{Dirichlet} (root-clamped) tree Laplacian. In particular each
block is positive-definite and the joint prior integrates to a
finite constant.
\end{theorem}
\begin{proof}
We argue axis by axis. Fix an axis $a$. The terms in
\eqref{eq:logprior-final} that depend on the $a$-axis entries of $x$
form a quadratic form
\(
Q_a(x) = (1/2\sigma_a^2)\sum_{v\neq\troot}\Delta_a(v)^2
       (+\, (1/2\sigma_\troot^2)(x^{(a)}_\troot - \loss^{(a)}_\troot)^2
       \text{ on the gain axis})\).
We must show $Q_a$ is positive-definite.

\textbf{Gain axis ($a=\rate$).} For every edge $\pa(v)\to v$, the
term $\Delta_\rate(v)^2 = (x^{(\rate)}_v - x^{(\rate)}_{\pa(v)})^2$
contributes $+1/\sigma_\rate^2$ to the $(v,v)$ and
$(\pa(v),\pa(v))$ diagonals and $-1/\sigma_\rate^2$ to the
off-diagonals $(v,\pa(v))$ and $(\pa(v), v)$.  Summing over edges
gives exactly the unsigned graph Laplacian
$L_\Tree = D - A$ scaled by $1/\sigma_\rate^2$, where $D$ is the
degree matrix and $A$ the adjacency matrix.  This $L_\Tree$ has
nullspace spanned by the all-ones vector $\mathbf{1}$ (Brownian
motion has a free additive constant).  The root anchor
$(1/\sigma_\troot^2)(x^{(\rate)}_\troot - \loss^{(\rate)}_\troot)^2$
contributes $+1/\sigma_\troot^2$ to the $(\troot,\troot)$ diagonal,
which removes that nullspace direction: any non-zero
$y \in \R^N$ either has $y \not\perp \mathbf{1}$ (in which case
$L_\Tree$ already gives strictly positive contribution) or
$y = c\mathbf{1}$ with $c\neq 0$ (in which case the root-anchor term
gives $c^2/\sigma_\troot^2 > 0$). Hence $\Lambda_\rate \succ 0$.

\textbf{Duplication and length axes ($a\in\{\dup, \edgelen\}$).}
On these axes the root entry is not in $x$ but is held fixed at the
prior centre $\loss^{(a)}_\troot$. The edge-increment terms involving
the root,
$(x^{(a)}_{v} - \loss^{(a)}_\troot)^2$ for $v \in \ch(\troot)$, then
contribute $+1/\sigma_a^2$ to the $(v,v)$ diagonal but no
off-diagonal. This is the standard Dirichlet condition: the
``root-clamped'' tree Laplacian $L^\partial_\Tree$ obtained from
$L_\Tree$ by deleting the root row/column has full rank (a connected
graph minus one boundary node has full-rank Laplacian).  Hence
$\Lambda_a \succ 0$.

The three axes do not interact in \eqref{eq:logprior-final}, so the
full precision is block-diagonal and positive-definite. The joint
density is a non-degenerate Gaussian and integrates to
$(2\pi)^{(3N-2)/2}\det(\Sigma)^{1/2}$, finite. $\square$
\end{proof}

\begin{remark}[Why the root anchor on $\rate$ is necessary]\label{rem:anchor}
Without the root anchor, $\Lambda_\rate$ is the bare tree Laplacian
$L_\Tree$, which is rank-deficient: it admits a one-dimensional
nullspace corresponding to a common additive shift of all
$\log\rate_v$. Biologically that is the ``slide the entire gain rate
up and down'' direction --- the data \emph{can} pin it (a high gain
rate predicts many copies; a low gain rate predicts few), but the
\emph{prior} cannot, and an improper prior on a direction the data
also struggles to constrain would propagate to a non-identifiable
posterior. The wide anchor $\sigma_\troot = 5$ contributes very
little except in the marginal axis the data cannot fix; it is the
minimum surgery to make the prior proper.
\end{remark}

\subsection{Coercivity of the MAP objective}
\label{sec:coercivity}

The point of using the prior at all is that the GLD likelihood
$\Ll(x)$ is not bounded above on $\R^{3N-2}$: at certain boundary
configurations the supremum is finite but unattained on $\R^{3N-2}$
(see the companion note \texttt{boundary\_identifiability.tex} for
the Pólya--Poisson asymptote analysis). The Brownian prior penalises
those configurations, and combined with the likelihood gives a
coercive objective.

\begin{theorem}[Coercivity of the Brownian-MAP objective]
\label{thm:coercive}
For any finite $\sigma_a, \sigma_\troot > 0$, the objective
\begin{equation}
\Ll^{\mathrm{MAP, B}}(x)
\;:=\;
\Ll(x) \;-\; F\log\bigl(1 - L(0;x)\bigr) \;+\; \log p_{\mathrm{B}}(x)
\label{eq:map-obj}
\end{equation}
is coercive on $\R^{3N-2}$: $\Ll^{\mathrm{MAP, B}}(x_n) \to -\infty$
whenever $\|x_n\| \to \infty$. Its supremum is therefore attained at
an interior point.
\end{theorem}
\begin{proof}
The likelihood term $\Ll(x) = \sum_f \log L_f(x)$ is non-positive
($L_f \in (0,1]$ for every family). The conditional-observability
correction $-F\log(1-L(0;x))$ is positive but bounded above by
$F\log\bigl(\liminf_{x\in\R^{3N-2}}(1 - L(0;x))\bigr)^{-1}$, which is
finite as long as the family-size process is non-degenerate (it is,
on any tree). Hence the first two terms are bounded above.

The third term $\log p_{\mathrm{B}}(x)$ is a quadratic form
$-\tfrac{1}{2}(x-\mu)^\top \Sigma^{-1}(x-\mu) + \mathrm{const}$, where
$\Sigma^{-1}$ is positive-definite by Theorem~\ref{thm:proper} and
$\mu$ is the prior mean. Hence
$\log p_{\mathrm{B}}(x) \le -\tfrac{1}{2}\lambda_{\min}(\Sigma^{-1})
\|x - \mu\|^2 + \mathrm{const}$,
where $\lambda_{\min}(\Sigma^{-1}) > 0$.  As $\|x\|\to\infty$ this
tends to $-\infty$ at quadratic rate, dominating any finite bounded
likelihood term, so $\Ll^{\mathrm{MAP, B}}(x) \to -\infty$.

Coercivity together with continuity of $\Ll^{\mathrm{MAP,B}}$
(continuity of $\Ll$ is standard from the smoothness of the survival
map of CM09; the prior is a smooth quadratic; $L(0;\cdot)$ is smooth
by the Pn-tensor recursion of Cs\H{u}r\"os 2026 SI) implies that
every super-level set $\{x : \Ll^{\mathrm{MAP, B}}(x) \ge c\}$ is
non-empty for some $c$ and bounded by coercivity, hence compact, on
which a continuous function attains its supremum. $\square$
\end{proof}

% =================================================================
\section{Default hyperparameters and biological calibration}
% =================================================================

\begin{table}[h]
\centering
\renewcommand{\arraystretch}{1.15}
\begin{tabular}{lll}
\toprule
parameter & default & biological meaning \\
\midrule
$\sigma_\rate$ & $1.0$ & one $e$-fold per edge in gain shape \\
$\sigma_\dup$ & $1.0$ & one logit-unit per edge in duplication \\
$\sigma_\edgelen$ & $1.0$ & one $e$-fold per edge in branch length \\
$\sigma_\troot$ & $5.0$ & deliberately wide; identifiability only \\
$\loss^{(\rate)}_\troot$ & $\log 0.1$ & matches \texttt{--mu-gain} \\
$\loss^{(\dup)}_\troot$ & $\logit(0.5) = 0$ & matches \texttt{--mu-dup} \\
$\loss^{(\edgelen)}_\troot$ & $\log 1 = 0$ & matches \texttt{--mu-length} \\
\bottomrule
\end{tabular}
\caption{Default hyperparameters. ``One $e$-fold per edge'' means
the 68\,\% Brownian-walk prior interval after one edge step on a
given axis is roughly $\pm 1$ on the log-scale, i.e.\ multiplicative
factors in $[e^{-1}, e^{1}] \approx [0.37, 2.72]$. Across a
20-edge-deep tree, the cumulative prior credible interval on
log-rate is then roughly $\pm \sqrt{20} \approx \pm 4.5$, i.e.\
$\sim 90\times$ either direction --- wide enough not to override
data, tight enough to forbid pathological boundary excursions of
$\sim 10^{10}\times$.}
\end{table}

\paragraph{What ``$\sigma_a = 1$'' is asking biologically.}
With $\sigma_a = 1$, two sister branches off a recent common
ancestor are expected to have log-rates within $\pm 2$ ($\sim$95\,\%
prior CI under a one-step Brownian) --- i.e.\ their rates within a
factor of $\sim 7$. That matches the kind of within-clade rate
heterogeneity reported in molecular-clock studies (where the
relaxed-clock literature finds factors of 3--10 across closely
related taxa). Tightening to $\sigma_a = 0.5$ enforces visibly more
smoothing; loosening to $\sigma_a \gg 1$ recovers the
independent-prior limit.

\paragraph{CLI coordinate convention.}
The CLI flag \texttt{--mu-dup} in \texttt{validation/map.py} is
specified by the user in \emph{log-rate units} (default
$\log 0.5 \approx -0.693$, i.e.\ ``prior centre at $\dup = 0.5$''),
not in the optimiser's logit-coord units. When the sub-critical
constraint is active the optimiser acts on
$\theta_{\dup,v} = \logit(\dup_v / \mathrm{Q})$, so the prior centre
$\loss^{(\dup)}_\troot$ appearing in \eqref{eq:logprior-final} must
also be in logit-coords. The translation happens once at the
boundary of \texttt{fit\_dataset} via \texttt{\_mu\_dup\_for\_optimizer};
each summary JSON records both the user-passed log-rate value and
the translated logit-coord value (\texttt{mu\_dup\_in\_optimizer\_coord}).

% =================================================================
\section{Implementation}
% =================================================================

\begin{itemize}
\item \textbf{Native C:} \texttt{native/src/recount\_brownian\_prior.c}.
  Public symbol \texttt{recount\_brownian\_prior\_and\_grad}, single
  $\mathcal{O}(N)$ loop over edges, no allocation, no threading
  (work is too small to benefit from \texttt{dispatch\_apply}).
\item \textbf{ctypes wrapper:}
  \texttt{recount.native\_backend.brownian\_prior\_native}, accepting
  the parent array, root index, packed $x$, the four sigmas, and the
  three prior centres; returning $(\log p_{\mathrm{B}}, \nabla\log p_{\mathrm{B}})$.
\item \textbf{Python reference:}
  \texttt{validation.\_shared.\_brownian\_prior\_and\_grad\_python}
  uses \texttt{numpy.add.at} for the scatter-add and is kept for
  finite-difference cross-checking in
  \texttt{tests/test\_brownian\_prior.py}.
\item \textbf{Integration with the existing MAP factory:}
  \texttt{make\_objgrad\_map\_brownian} in
  \texttt{validation/\_shared.py} subtracts $\partial\log p_{\mathrm{B}}/\partial x$
  from the negative-LL gradient produced by
  \texttt{\_gradient\_raw\_native}; the LL chain rule is unchanged.
\item \textbf{Validation:} analytical gradient matches central
  finite differences to relative error $\le 10^{-6}$ on 30 random
  optimiser positions on \texttt{dpann80} ($N=159$); the
  full-objective gradient through
  \texttt{make\_objgrad\_map\_brownian} matches FD to the same
  tolerance. See \texttt{tests/test\_brownian\_prior.py} for the
  four-test suite.
\end{itemize}

% =================================================================
\section{Out of scope for v1, deferred to v2}
\label{sec:future}
% =================================================================

\begin{enumerate}[leftmargin=*]
\item \textbf{TKP-classic per-edge $\edgelen_v$ weighting}
  (\eqref{eq:tkp-classic}, the proper diffusion form). Requires extra
  chain-rule terms since $\edgelen_v$ is itself in the optimiser.
  \emph{Consequence of the v1 simplification}: very long branches are
  penalised as tightly as very short ones for a 1-unit log-rate
  change. On real phylogenies with heterogeneous branch lengths, the
  v1 prior systematically over-smooths short branches and
  under-smooths long branches relative to the proper TKP form.
\item \textbf{Empirical Bayes / REML estimation of the $\sigma_a$.}
  V1 fixes the three sigmas from the CLI; an outer EM or REML loop
  is a v2 enhancement.
\item \textbf{Composition with the $K$-category LogisticShift mixture}
  (\texttt{logistic\_shift\_gradient.tex}). The current mixture fit
  applies an independent Gaussian prior on the base rates; replacing
  it with the Brownian prior would compose family-level heterogeneity
  with branch-level smoothing.
\item \textbf{Per-branch discrete rate-class HMM.} A separate model
  class with its own forward--backward, beyond the scope of the
  Brownian-prior framework.
\end{enumerate}

% =================================================================
\section{Related work}
% =================================================================

\begin{itemize}
\item \textbf{Thorne, Kishino \& Painter} (1998),
  \emph{Estimating the rate of evolution of the rate of molecular
  evolution}, MBE 15:1647--1657. The original autocorrelated-rate
  prior; we adapt the structure unchanged.
\item \textbf{Drummond, Ho, Phillips \& Rambaut} (2006),
  \emph{Relaxed phylogenetics and dating with confidence}, PLoS
  Biology 4:e88. The uncorrelated-lognormal alternative
  ($\log r_v$ iid from a global distribution, ignoring tree
  structure). The recount independent prior is closer to this
  unstructured form; the Brownian prior described here is its TKP
  autocorrelated counterpart.
\item \textbf{Cs\H{u}r\"os \& Mikl\'os} (2009), \emph{Streamlining and
  large ancestral genomes in Archaea inferred with a phylogenetic
  birth-and-death model}, MBE 26:2087--2095. The GLD model whose
  per-node rates we are placing a prior on. Cs\H{u}r\"os' Java
  reference (\texttt{count.model.MLDistribution}) does
  plain maximum-likelihood optimisation of the rates: no Bayesian
  prior on the per-node $(\rate_v,\dup_v,\edgelen_v)$. The only
  ``prior'' in Cs\H{u}r\"os' code is the \emph{root prior} on the
  copy count at the root, which is part of the model itself, not a
  regulariser on the rate parameters; and the ``MAP'' language in
  Cs\H{u}r\"os 2026 PNAS refers to maximum-a-posteriori
  family-presence inference (a per-family binary classifier on the
  inferred copy counts at each internal node), not MAP rate
  estimation. The independent per-node log-Normal prior in
  \texttt{validation/map.py} as the \texttt{--prior independent}
  mode, and the Brownian prior described here as
  \texttt{--prior brownian}, are both additions by the present
  codebase.
\item \textbf{Cs\H{u}r\"os} (2021), arXiv:2107.11440. Analytical
  gradient of the GLD likelihood via the conserved-copy
  decomposition. Provides the per-rate gradients used in our
  $\Ll(x)$ pipeline.
\end{itemize}

\end{document}
